\documentclass{beamer}
\usepackage[latin1]{inputenc}
\usepackage{graphicx}
\usepackage{amssymb}
\usepackage{amsmath}
\usepackage{listings}
\usepackage{hyperref}

\usetheme{default}

\title[titre1]{TITRE de ma PRESENTATION}
\author[A. Sanchez]{Arthur Sanchez}
\date[]{Juillet 2022\\Paris}

\begin{document}

\begin{frame}[plain]
\titlepage
\end{frame}

\begin{frame}
\frametitle{Plan}
  \tableofcontents
\end{frame}

\section{Premi�re section}

\begin{frame}
  \frametitle{Titre premi�re diapo}
  \framesubtitle{Sous-titre premi�re diapo}
  Une liste d'items...
  \begin{itemize}
  \item item 1
  \item \textbf{item 2} %texte en gras
  \item \textit{item 3} %texte en italique
  \end{itemize}
\end{frame}

\begin{frame}
  \frametitle{Titre deuxi�me diapo (sans sous titre)}
  Une liste d'items qui s'affiche progressivement
  \begin{itemize}
  \onslide<1-3> \item item 1
  \onslide<2-3>\item \textbf{item 2} %texte en gras
  \onslide<3-3>\item \textit{item 3} %texte en italique
  \end{itemize}
\end{frame}

\begin{frame}
  \frametitle{Enumerate plut�t que itemize pour des num�ros..}
  Une liste d'items num�rot�s...
  \begin{enumerate}
  \item item 1
  \item \textbf{item 2} %texte en gras
  \item \textit{item 3} %texte en italique
  \end{enumerate}
\end{frame}

\begin{frame}
  \frametitle{Un exemple de tableau...}
  \begin{table}
    \begin{tabular}{| c || c | c |}
      \hline
      No. & Name & Age \\
      \hline \hline
      1 & John T & 24 \\
      2 & Norman P & 8 \\
      3 & Alex K & 14 \\ 
      \hline
    \end{tabular}
    \caption{Name and age of students}
  \end{table}
\end{frame}

\section{Images et sch�mas}

\begin{frame}
  \frametitle{Un exemple d'image}
  \begin{center}
    \includegraphics[width=4cm]{image.jpg}
  \end{center}
\end{frame}

\begin{frame}
  \frametitle{Un exemple de sch�ma/dessin}
  % sauvegarder dans inkscape le document en svg puis sauvegarder comme pdf 
  % ne pas oublier de sauvegarder en svg
  \begin{center}
    \includegraphics[width=4cm]{dessin.pdf} 
  \end{center}
\end{frame}

\section{Code et formules math�matiques}

\begin{frame}[fragile]
  \frametitle{Un premier exemple de code...}
  \begin{small}
  \begin{lstlisting}[numbers=left,frame=trBL]
# Python Program to find the area of triangle

a = 5
b = 6
c = 7

# Uncomment below to take inputs from the user
# a = float(input('Enter first side: '))
# b = float(input('Enter second side: '))
# c = float(input('Enter third side: '))

# calculate the semi-perimeter
s = (a + b + c) / 2

# calculate the area
area = (s*(s-a)*(s-b)*(s-c)) ** 0.5
print('The area of the triangle is %0.2f' %area)
\end{lstlisting}
\end{small}
\end{frame}

\begin{frame}[fragile]
  \frametitle{Un deuxi�me exemple de code...}
    \begin{small}
      \lstinputlisting[numbers=left,language=python]{area-triangle.py}
    \end{small}
\end{frame}

\begin{frame}[fragile]
  \frametitle{Un dernier exemple de code...}
    \lstinputlisting[language=python, linerange={3-5,15-17}]{area-triangle.py}
\end{frame}


\begin{frame}
  \frametitle{Un exemple de formule math�matiques...}
  $\forall x \in \mathbf{R}: \qquad x^{2} \geq 0$
  \ \\ \ \\ \ \\ % ins�rer quelques lignes blanches
  Quelques exemples de formules \href{https://www.math-linux.com/latex-4/faq/latex-faq/article/latex-derivee-limite-somme-produit-et-integrale}{ICI}
\end{frame}

\section{Quelques autres fonctionnalit�s utiles}

\begin{frame}
  \frametitle{Pour diviser la diapositive en 2 colonnes}
  \begin{columns}
    % Column 1
    \begin{column}{0.5\textwidth}
      \begin{itemize}
      \item blab bla bla
      \item blab bla bla
      \item blab bla bla
      \end{itemize}
    \end{column}

    % Column 2
    \begin{column}{0.5\textwidth}
      \includegraphics[scale=0.05]{image.jpg}
    \end{column}
  \end{columns}
\end{frame}

\begin{frame}
  \frametitle{Misc...}
  Pour aller � la ligne...\\
  Et avec une ligne de plus \\ \ \\
  et pour �crire \textcolor{blue}{en couleur} %blue green yellow orange red violet
\end{frame}


\end{document}


